\section{Introduction} \label{sec:intro}

Fast radio bursts (FRBs) are short-duration ($\lesssim 100$\,ms), highly energetic ($\gtrsim 10^{40}$\,erg\,s$^{-1}$) transients that have been detected between 100\,MHz and 8\,GHz~\cite{cordesFastRadioBursts2019, petroffFastRadioBursts2019}. 
Prior to April 2020, all FRBs were extragalactic and most were at cosmological distances ($z\geq0.1$).
The discovery of SGR\,1935+2154 by both STARE2 and the Canadian Hydrogen Intensity Mapping Experiment (CHIME/FRB) made a dramatic connection between extragalactic FRBs and Galactic magnetars~\cite{kirstenDetectionTwoBright2021,bochenekSTARE2DetectingFast2020}.
While FRB\,200428 from SGR\,1935+2154 was less luminous than all known extragalactic FRBs, its proximity made it the highest fluence radio burst ever detected ($\sim$\,MJy\,ms at 1\,GHz) and enabled detection of the first prompt multi-wavelength emission~\cite{mereghettiINTEGRALDiscoveryBurst2020}.

In the years since the FRB-like emission from SGR\,1935+2154, the link between young, energetic magnetars and the extragalactic FRB phenomenon has been muddied.
FRBs have been found in globular clusters~\cite{bhardwajNearbyRepeatingFast2021}, at the outskirts of elliptical galaxies~\cite{shahRepeatingFastRadio2025, eftekhariMassiveQuiescentElliptical2025}, and may be underrepresented in low-mass star forming galaxies and the sites of recent core-collapse supernovae~\cite{sharmaPreferentialOccurrenceFast2024}.
In other words, if all FRBs are magnetars, they are not all young remnants of CCSNe embedded in star-forming region, like SGR\,1935+2154.
These facts motivate a blind search survey, rather than monitoring known Galactic magnetars.
The discovery of other new Galactic phenomena such as long-period radio transients~\cite{hurley-walkerRadioTransientUnusually2022, calebEmissionstateswitchingRadioTransient2024, rodriguezSpectroscopicDetection29hour2025} further motivates blind searches in new regions of parameter space. 

The Galactic Radio Explorer is an international collaboration aimed at discovering exceptionally bright bursts in the radio sky.
Building off the success of STARE2~\cite{bochenekSTARE2DetectingFast2020}, we are currently searching for FRBs emitted by Galactic sources as well as nearby galaxies.
We search down to a high time resolution of $32.768\,\mu\text{s}$, saving raw $8.192\,\mu\text{s}$ voltage data to enable more in-depth analysis for interesting single-pulse candidates.
Additionally, the instrument has been designed to be easy to build and reproduce, allowing collaborators to quickly bring up a functioning station.
The goal as mentioned in our original white paper~\cite{connorGalacticRadioExplorer2021a} is to eventually have 4$\pi$ steradian coverage with increased exposure to the galactic plane to improve sensitivity to galactic sources.
As more sensitive aperture arrays for fast transient discovery come online, such as BURSTT~\cite{linBURSTTBustlingUniverse2022} in Taiwan and similar efforts in the US, Australia, and Chile, the niche of GReX will be all-sky 
sensitivity to ultra-narrow, bright radio bursts at 1.4\,GHz.
Each GReX terminal is fully self-contained and low-cost, serving as a valuable pedagogical tool for the observatory or university operating it locally.

In this chapter, we describe the GReX instrument, and the current network of terminals around the world. We also place new upper-limits on the rate of ultra-bright FRBs based on non-detection.
